%% kapa: exam=1 round=20 year=2009 period=2 form=A subject=accounting subject_ko=회계학 seq=25 number=25
%% source: Q-Net 감정평가사 기출문제 (cst00309) artlSeq=5024178 file=extracted/1008/5024178/2009년 제20회 감정평가사 1차시험 기출문제/2교시.hwp
%% conversion: kordoc 4.24.51 -> tools/gen1.py (auto); review=proofread-sonnet; flags=none
\begin{question}{25}
다음은 (주)한국이 20×8년과 20×9년 중 단기매매목적으로 거래한 사채와 관련된 내용이다.
\begin{framedtext}
ㆍ20×8년 10월 1일에 (주)서울이 발행한 사채(액면금액 ￦1,000,000, 이자율 연 10\%, 이자는 매년 말 후급, 20×8년 1월 1일 발행, 20×9년 12월 31일 원금일시상환)를 발생이자를 포함하여 ￦1,050,000에 취득하였다. \par
ㆍ20×8년 말에 현금이자를 수령하였으며 기말현재 상기사채와 관련하여 형성된 시장수익률은 연 10\%이다. \par
ㆍ20×9년 7월 1일에 상기사채 전액을 발생이자를 포함하여 ￦1,020,000에 처분하였다. \par
\end{framedtext}
위의 거래와 관련하여 (주)한국이 20×8년도에 인식할 단기매매증권평가손익과 20×9년도에 인식할 단기매매증권처분손익은 각각 얼마인가? \par

{\small\centering
\begin{tabular}{@{}ccc@{}}
 & 단기매매증권평가손익 & 단기매매증권처분손익 \\
\cline{2-2}\cline{3-3}
① & ￦50,000 손실 & ￦5,000 이익 \\
② & ￦25,000 이익 & ￦5,000 이익 \\
③ & ￦25,000 이익 & ￦30,000 손실 \\
④ & ￦50,000 손실 & ￦30,000 손실 \\
⑤ & ￦50,000 손실 & ￦20,000 이익 \\
\end{tabular}\par}
\end{question}
\answer{3}
